\chapter{Conventions}

\section*{Typefaces}

Latin\-/script Interturkic text is set in a serif typeface,
e.g.\ \latntxt{baş}, \latntxt{fikr}, \latntxt{tığ}.
Arabic\-/script Interturkic text is set in a Naskh\-/like typeface,
e.g.\ \arabtxt{باش}, \arabtxt{فکر}, \arabtxt{تیغ}.
When Latin\-/script text from any other language\---including Interturkic\---is
embedded within English\-/language prose, it is italicised,
e.g.\ \foreign{\latntxt{baş}}, \foreign{\latntxt{fikr}}, \foreign{\latntxt
{tığ}}.
Glossing abbreviations are set in small caps,
e.g.\ \Fpl, \Abl, \Pst.

\section*{Transcriptions}

All phonological transcriptions are rendered with symbols from the International
Phonetic Alphabet~(IPA).

Variable segments are represented by uppercase letters.
`A' represents the set \phonem{ɑ, ɛ}.
`D' represents the set \phonem{d, t}.
`G' represents the set \phonem{g, ɣ, k, q}.
`I' represents the set \phonem{ɯ, i, u, y}.
`K' represents the set \phonem{k, q}.
This is the case in both phonological transcriptions and in Latin\-/script
Interturkic text,
e.g.\ \foreign{\latntxt{\morphophon{\-/nIñ}}} \phonem{\morphophon{\-/nIŋ}}
\inlinegls{\Gen}.
